% ============================================================
% CHƯƠNG 7: KẾT LUẬN
% PHÂN BỔ: ~2 trang
% ============================================================
\chapter{Kết luận}
\label{chap:ket-luan}

% ============================================================
\section{Tóm tắt đóng góp Phase 1}
\label{sec:ket-luan-dong-gop}

Đồ án chuyên ngành \textbf{ELDAS} đặt mục tiêu xây dựng một nền tảng
mô phỏng và lập lịch tác vụ AI trong datacenter dựa trên học tăng
cường đa mục tiêu (MORL). Báo cáo này trình bày kết quả Phase 1 ---
giai đoạn xây dựng nền tảng kỹ thuật và thiết lập đường tham
chiếu --- với sáu đóng góp có thể kiểm chứng độc lập:

\begin{enumerate}
  \item \textbf{Kiến trúc bốn lớp với contract rõ ràng}
        (Chương~\ref{chap:kien-truc}). Bốn lớp Docker / CloudSim
        Plus / Py4J / MO-Gymnasium được tách rời tới mức có thể
        thay engine mô phỏng hoặc đổi thuật toán RL mà không cần
        chỉnh các lớp còn lại --- đây cũng là điều kiện tiên quyết
        cho Phase 2 và các hướng nghiên cứu mở rộng.
  \item \textbf{Engine mô phỏng có theo dõi GPU}
        (§\ref{ssec:impl-factory}). CloudSim Plus 8.5.7 nguyên bản
        chưa có khái niệm GPU; ELDAS bổ sung lớp \texttt{GpuState}
        thông qua \texttt{IdentityHashMap} kèm công thức năng lượng
        tuyến tính riêng cho GPU (\eqref{eq:gpu-power}), tránh phải
        fork thư viện gốc.
  \item \textbf{Tích hợp dữ liệu workload thực} từ Alibaba PAI
        Cluster Trace v2023. Parser chịu lỗi
        \texttt{AlibabaTraceReader} kết hợp bộ lọc
        \texttt{ScenarioFilter} cho ba kịch bản tải Low/High/Burst
        (§\ref{ssec:impl-trace}, §\ref{ssec:impl-baselines}).
  \item \textbf{Cầu nối ngôn ngữ tốc độ cao}
        (§\ref{sec:layer-py4j}). Py4J được lựa chọn thay vì gRPC
        hay REST nhờ độ trễ chỉ $\sim 0{,}1$\,ms; vận hành thực tế
        cho phép $\sim 970$ step/s qua bốn lớp
        (Bảng~\ref{tab:smoke-results}).
  \item \textbf{Vector reward kết hợp action masking}
        (§\ref{sec:layer-mogym}). Môi trường tuân thủ chuẩn
        MO-Gymnasium 1.0+, giữ reward ở dạng vector hai chiều; cơ
        chế \texttt{action\_masks()} của MaskablePPO bảo đảm chính
        sách không sinh ra hành động vi phạm \eqref{eq:capacity}.
  \item \textbf{Khả năng tái lập đến từng bit}
        (§\ref{ssec:impl-build}, §\ref{ssec:res-repro}). Docker
        Compose, version pinning và \texttt{RANDOM\_SEED} dùng
        chung cho cả Java lẫn Python. Bằng chứng: \texttt{diff}
        hai run cùng seed trả về output rỗng (Hình~\ref{fig:repro}).
\end{enumerate}

Kết quả thực nghiệm ở Chương~\ref{chap:ket-qua} đã đối chiếu thành
công cả sáu mục tiêu MT1--MT6 (Bảng~\ref{tab:mt-checklist}). Sau
khi mở rộng phạm vi sang Phase 1.8 (state-machine năng lượng cùng
ba baseline cổ điển bổ sung), \textbf{năm baseline} --- FirstFit,
BestFit, RoundRobin, K8s LRP và Random --- đã đặt sẵn đường tham
chiếu Pareto cho Phase 2. Trên kịch bản HIGH, chênh lệch năng lượng
giữa BestFit (pack-tight, $21\,817$\,kWh) và RoundRobin (spread,
$27\,369$\,kWh) lên đến $\mathbf{20{,}3\%}$; SLA reward đi theo
chiều ngược lại, K8s LRP tốt hơn BestFit khoảng $3{,}7$ lần. Đây
là một Pareto front \emph{không tầm thường} với ba cụm rõ rệt:
agent MORL phải cải thiện đồng thời trên cả hai trục, chứ không
chỉ một trục như giả định ban đầu của báo cáo.

% ============================================================
\section{Bài học kinh nghiệm}
\label{sec:bai-hoc}

Bốn bài học kỹ thuật nổi bật mà bản thân tác giả rút ra được trong
quá trình hiện thực:

\paragraph{Giữ chặt biên giới giữa hai ngôn ngữ.}
Khi tích hợp Java và Python, cám dỗ thường gặp là phơi luôn các cấu
trúc dữ liệu nội bộ ra cả hai phía. Quyết định giữ
\texttt{GatewayEntryPoint} chỉ phơi các \emph{record nguyên thuỷ}
(Bảng~\ref{tab:gw-api}) đã chứng tỏ giá trị của mình: khi cần đổi
backend, biên giới giữa hai ngôn ngữ chỉ còn là chín phương thức gọn,
chứ không phải hàng trăm field rải rác.

\paragraph{Cơ chế pause/resume là phần tinh tế nhất.}
Cặp \texttt{SynchronousQueue} cho phép tránh hoàn toàn việc dùng
khoá tường minh mà vẫn đồng bộ chính xác giữa hai luồng. Đây là một
mẫu thiết kế chuẩn cho việc tích hợp giữa một \emph{event-driven
simulator} và một thuật toán RL on-policy, và rất đáng để áp dụng
lại ở các dự án tương lai.

\paragraph{Đôi khi theo dõi thủ công lại tốt hơn ``dùng đúng'' API.}
Việc cố ý \emph{không} gọi \texttt{simulation.start()} ở Phase 1
(§\ref{ssec:impl-sim-manager}) là một quyết định có chủ đích, được
ghim chặt bằng comment trong source. Bài học chung: khi một engine bên thứ ba có hành
vi tiềm tàng nguy hiểm, đi đường vòng \emph{có kiểm soát} an toàn hơn
là cố ``vá cho hết warning''.

\paragraph{Khả năng tái lập là kết quả của kỷ luật, không phải may mắn.}
Mọi seed, mọi version và mọi đường dẫn đều được pin chặt; trace có
hash MD5; hai run với cùng tham số thực sự cho output trùng đến từng
byte. Khi công bố kết quả nghiên cứu, đây là cam kết tối thiểu cần
phải có, không thể thoả hiệp.

% ============================================================
\section{Lời kết}
\label{sec:loi-ket}

Bài toán lập lịch tác vụ AI trong datacenter dưới các ràng buộc đa
mục tiêu đang ngày càng tốn kém với ngành công nghiệp tính toán
hiện đại: chỉ một phần trăm năng lượng tiết kiệm được trên một
cluster A100 quy mô vừa cũng đã tương đương với hàng chục nghìn USD
chi phí điện mỗi năm, cùng hàng tấn $\mathrm{CO_2}$ tránh được phát
thải. Phase 1 của ELDAS đã đặt nền móng để câu hỏi này có thể được
đo lường, lặp lại và so sánh một cách công bằng trong môi trường
nghiên cứu. Phase 2 sẽ tiếp tục trả lời câu hỏi định lượng kế tiếp:
liệu MORL có thực sự vượt được Kubernetes LRP --- một baseline đã
được vận hành trong sản xuất suốt nhiều năm --- hay không? Dù câu
trả lời là ``có'' hay ``không'', kết quả này đều có giá trị tham
khảo cho cộng đồng nghiên cứu trong nước và quốc tế.

\bigskip
\noindent
Hệ thống ELDAS được công bố dưới dạng mã nguồn mở, kèm cấu hình
Docker Compose đầy đủ, với mong muốn trở thành điểm khởi đầu cho
các đề tài kế tiếp về lập lịch xanh, học tăng cường đa mục tiêu và
tối ưu hệ thống phân tán --- những hướng nghiên cứu thiết yếu trong
lộ trình Net-Zero $2050$ của Việt Nam.

\begin{flushright}
\textit{TP. Hồ Chí Minh, tháng 05 năm 2026}\\
\textbf{Vũ Trần Duy Nguyên}
\end{flushright}
